% !TeX spellcheck = ru_RU
% !TEX root = report.tex

\section*{Введение}

Сложные структурированные объекты~--- сцены, молекулы, схемы~--- удобно описывать как множества элементов, обладающих некоторыми свойствами и связанных некоторыми отношениями.
Традиционный язык таких описаний~--- формулы исчисления предикатов.
Задачи распознавания в такой постановке NP-трудны~\cite{Kosovskaya2007}, поэтому Т.~М.~Косовская с соавторами предложили выделять общие для класса объектов подформулы-шаблоны; в основе подхода лежит понятие изоморфизма элементарных конъюнкций через шаблон с переменными~\cite{ZhouKosovskaya2024, KosovskayaZhou2024}.

У формул исчисления предикатов есть неудобство: каждое отношение имеет фиксированную арность и упорядоченные аргументы, так что утверждение \enquote{пятеро людей~--- братья} требует 20 атомарных формул.
Поэтому в работе~\cite{KosovskayaHypergraphs2025} предложены \emph{гиперграфы отношений}, где коммутативное отношение любой арности задаётся одним гиперребром, а в статье~\cite{KosovskayaIsom}~--- алгоритмы проверки их изоморфизма \ISOMone{} и \ISOMn{}.
Алгоритмы описаны словесно и блок-схемами и проиллюстрированы примерами; программной реализации и экспериментальной оценки у них пока нет.

Чтобы пользоваться алгоритмами, их нужно реализовать, проверить реализацию на независимом эталоне и понять, как время работы зависит от вида входных данных.
Важен и выбор языка: на Python алгоритм удобно записывать близко к тексту статьи и исследовать, а на \Cpp{} он работает быстрее.
Интересно, насколько велика эта разница, если обе реализации выполняют в точности один и тот же поиск.

В работе алгоритмы статьи реализованы на Python и на \Cpp{} так, что обе реализации строят одинаковые деревья поиска.
Корректность проверена сверкой с полным перебором и сверкой реализаций между собой, а производительность сравнена на десяти наборах задач, от примеров из статьи до гиперграфов с сотнями вершин.
